\chapter{Numeri naturali e principio di induzione}\label{cap:induzione}

Finora abbiamo usato i numeri naturali negli esempi come oggetti già noti. In questo capitolo vediamo come la teoria degli
insiemi li fa nascere: da un insieme infinito si ricava un insieme ordinato con le proprietà di $\mbN$, e da queste proprietà
segue lo strumento più importante per dimostrare enunciati «per ogni $n$»: il \textbf{principio di induzione}. Lo useremo subito
per dimostrare le regole delle potenze nei semigruppi, rimaste in sospeso nella Parte~III.

Le operazioni $+$ e $\cdot$ su $\mbN$ e su $\mbZ$, con le loro proprietà, continuiamo a darle per note.

% ──────────────────────────────────────────────────────────────
\section{Insiemi infiniti}\label{sec:infiniti}

\dfn[equipotenti]{Insiemi equipotenti}{
    Due insiemi $a$ e $b$ si dicono \textbf{equipotenti} se esiste una funzione biettiva $f \colon a \to b$.
}

\begin{concetto}
\begin{Proposizione}{Proprietà dell'equipotenza}{equipotenza}
    \begin{prerequisiti}
        \dip{Definizione}{def:equipotenti}
        \dip{Teorema}{teo:biettiva-inversa}
        \dip{Corollario}{cor:comp-biettive}
    \end{prerequisiti}
    Siano $a$, $b$, $c$ insiemi. Allora $a$ è equipotente ad $a$; se $a$ è equipotente a $b$, allora $b$ è equipotente ad $a$; se $a$ è
    equipotente a $b$ e $b$ a $c$, allora $a$ è equipotente a $c$.
\end{Proposizione}
\begin{dimostrazione}
    L'identità $\mathrm{id}_a$ è biettiva. Se $f \colon a \to b$ è biettiva, ha un'inversa $f^{-1} \colon b \to a$, che è biettiva
    (Teorema~\ref{teo:biettiva-inversa}: la sua inversa è $f$). Se $f \colon a \to b$ e $g \colon b \to c$ sono biettive, $g \circ f$
    è biettiva (Corollario~\ref{cor:comp-biettive}).
\end{dimostrazione}
\end{concetto}

\begin{spiegazione}
    Due insiemi sono equipotenti quando i loro elementi si possono \textbf{accoppiare uno a uno} senza avanzi, come i posti e gli
    spettatori in un teatro tutto esaurito. Questo accoppiamento si può rovesciare (inversa) e concatenare (composta): per questo
    l'equipotenza si comporta come una relazione di equivalenza. Per gli insiemi finiti vorrà dire «avere lo stesso numero di
    elementi»; per quelli infiniti riserva qualche sorpresa.
\end{spiegazione}

\dfn[infinito]{Insieme infinito}{
    Un insieme si dice \textbf{infinito} se è equipotente a una sua parte propria.
}

La definizione sembra strana, ma coglie proprio ciò che distingue gli insiemi infiniti: un insieme finito non può essere
messo in corrispondenza biunivoca con una sua parte più piccola (lo vedremo nel Teorema~\ref{teo:in-non-infinito}), mentre
per un insieme infinito questo è possibile (Esempio~\ref{es:hilbert}). Che ne esista almeno uno non si può dimostrare dagli
assiomi visti finora: è un nuovo assioma.

\ass[cantor]{Dell'infinito (o di Cantor)}{
    Esiste un insieme infinito.
}

% ──────────────────────────────────────────────────────────────
\section{I numeri naturali}\label{sec:naturali}

\dfn[naturalmente-ordinato]{Insieme naturalmente ordinato}{
    Sia $(s, \leq)$ un insieme ordinato. Una parte $t \subseteq s$ si dice \textbf{superiormente limitata} se ha almeno un maggiorante.
    $(s, \leq)$ si dice \textbf{naturalmente ordinato} se è ben ordinato e ogni sua parte non vuota superiormente limitata ha massimo.
}

I due teoremi seguenti dicono che, a meno di isomorfismi, esiste uno e un solo insieme naturalmente ordinato senza maggioranti.

\begin{concetto}
\begin{Teorema}{Infinito e insiemi naturalmente ordinati}{infinito-naturale}
    \begin{prerequisiti}
        \dip{Definizione}{def:infinito}
        \dip{Definizione}{def:naturalmente-ordinato}
    \end{prerequisiti}
    \senzadim
    Esiste un insieme infinito se e solo se esiste un insieme naturalmente ordinato non superiormente limitato.
\end{Teorema}
\end{concetto}

\begin{concetto}
\begin{Teorema}{Unicità dei naturali}{unicita-naturali}
    \begin{prerequisiti}
        \dip{Definizione}{def:naturalmente-ordinato}
        \dip{Definizione}{def:crescente}
    \end{prerequisiti}
    \senzadim
    Tutti gli insiemi naturalmente ordinati non superiormente limitati sono isomorfi (come insiemi ordinati).
\end{Teorema}
\end{concetto}

\dfn[naturali]{I numeri naturali}{
    Per l'assioma di Cantor e il Teorema~\ref{teo:infinito-naturale} esiste un insieme naturalmente ordinato non superiormente
    limitato; ne scegliamo uno, lo indichiamo con $(\mbN, \leq)$ e lo chiamiamo \textbf{insieme dei numeri naturali}
    (per il Teorema~\ref{teo:unicita-naturali} la scelta è indifferente). Inoltre:
    \begin{itemize}
        \item $0 := \min(\mbN)$, $1 := \min(\mbN \setminus \{0\})$, e così via;
        \item per $n \in \mbN$, il \textbf{successore} di $n$ è $n + 1 := \min\{k \in \mbN \mid n < k\}$;
        \item per $n \in \mbN \setminus \{0\}$, il \textbf{predecessore} di $n$ è $n - 1 := \max\{k \in \mbN \mid k < n\}$;
        \item per $m \in \mbN$, poniamo $\mbN_m := \{n \in \mbN \mid m \leq n\}$ (i naturali da $m$ in poi).
    \end{itemize}
}

\begin{concetto}
\begin{Proposizione}{Successore e predecessore}{successore}
    \begin{prerequisiti}
        \dip{Definizione}{def:naturali}
        \dip{Definizione}{def:naturalmente-ordinato}
        \dip{Teorema}{teo:ben-ordinato-totale}
    \end{prerequisiti}
    Il successore e il predecessore esistono, e per ogni $n \in \mbN$:
    \begin{enumerate}
        \item $n < n + 1$ e non esiste $k$ con $n < k < n + 1$ (cioè $n + 1$ copre $n$);
        \item se $n \neq 0$: $n - 1 < n$, non esiste $k$ con $n - 1 < k < n$, e $(n - 1) + 1 = n$.
    \end{enumerate}
\end{Proposizione}
\begin{dimostrazione}
    Ricordiamo che $\mbN$ è totalmente ordinato (Teorema~\ref{teo:ben-ordinato-totale}).
    \begin{caso}{Esistenza}
        $\{k \mid n < k\}$ non è vuoto, perché altrimenti $n$ sarebbe un maggiorante di $\mbN$; essendo $\mbN$ ben ordinato, ha minimo.
        Se $n \neq 0$, $\{k \mid k < n\}$ contiene $0$ e ha il maggiorante $n$: per l'ordine naturale ha massimo.
    \end{caso}
    \begin{caso}{Punto 1}
        $n + 1$ sta in $\{k \mid n < k\}$, quindi $n < n + 1$; se fosse $n < k < n + 1$, $k$ sarebbe in quell'insieme e più piccolo del minimo.
    \end{caso}
    \begin{caso}{Punto 2}
        Allo stesso modo $n - 1 < n$ e nessun $k$ sta strettamente tra i due. Infine $n$ sta in $\{k \mid n - 1 < k\}$; se $j$ sta in
        quest'insieme e fosse $j < n$, allora $j \leq n - 1$ (massimo), contro $n - 1 < j$. Quindi $j \geq n$ per ogni tale $j$, e il
        minimo, cioè $(n - 1) + 1$, è $n$.
    \end{caso}
\end{dimostrazione}
\end{concetto}

\begin{spiegazione}
    In $\mbN$ ogni numero ha un \textbf{numero immediatamente dopo} e (tranne $0$) un \textbf{numero immediatamente prima}: è
    proprio ciò che manca, ad esempio, in $(\mbQ, \leq)$, dove tra due numeri ce n'è sempre un altro. Il diagramma di Hasse di
    $(\mbN, \leq)$ è una catena infinita verso l'alto, in cui ogni numero è coperto dal suo successore:
    \begin{center}
    \begin{tikzpicture}[hasse]
        \foreach \n in {0,...,5} \node[hv] (n\n) at (\n*1.3,0) {$\n$};
        \foreach \n [evaluate=\n as \m using int(\n+1)] in {0,...,4} \draw[hlato, -{Stealth[length=2mm]}] (n\n) -- (n\m);
        \node at (7.2,0) {$\cdots$};
        \draw[hlato, -{Stealth[length=2mm]}] (n5) -- (6.8,0);
        \node[font=\small] at (3.25,-0.65) {la freccia va da $n$ al suo successore $n + 1$ (disegniamo la catena in orizzontale)};
    \end{tikzpicture}
    \end{center}
\end{spiegazione}

\begin{esempio}{L'albergo di Hilbert}\label{es:hilbert}
    La funzione $n \in \mbN \mapsto n + 1 \in \mbN \setminus \{0\}$ è biettiva: ogni $m \neq 0$ è immagine del solo $m - 1$. Quindi
    $\mbN$ è equipotente alla sua parte propria $\mbN \setminus \{0\}$, ed è infinito. È il famoso «albergo di Hilbert»: un albergo
    con infinite stanze tutte occupate può accogliere un nuovo ospite, se ogni cliente si sposta dalla stanza $n$ alla stanza $n + 1$
    e lascia libera la stanza $0$.
\end{esempio}

% ──────────────────────────────────────────────────────────────
\section{Il principio di induzione}\label{sec:induzione}

\begin{concetto}
\begin{Teorema}{Principio di induzione di prima forma}{induzione-prima}
    \begin{prerequisiti}
        \dip{Definizione}{def:naturali}
        \dip{Proposizione}{prop:successore}
        \dip{Definizione}{def:ben-ordinato}
    \end{prerequisiti}
    Sia $x$ una parte non vuota di $\mbN$ tale che $\forall n \in \mbN\,(n \in x \rightarrow n + 1 \in x)$. Allora $x = \mbN_{\min(x)}$. In formule:
    \[ \forall x \in \mcP(\mbN) \setminus \{\varnothing\}\ \Big(\forall n \in \mbN\,\big(n \in x \rightarrow n + 1 \in x\big) \rightarrow x = \mbN_{\min(x)}\Big). \]
\end{Teorema}
\begin{dimostrazione}
    Sia $m = \min(x)$ (esiste perché $\mbN$ è ben ordinato). Ogni elemento di $x$ è $\geq m$, quindi $x \subseteq \mbN_m$.
    Per assurdo sia $x \neq \mbN_m$: allora $y = \mbN_m \setminus x$ non è vuoto e ha un minimo $n$.
    \begin{passi}
        \item $n \neq m$, perché $m \in x$ e $n \notin x$; quindi $m < n$, e in particolare $n \neq 0$.
        \item $n - 1 \geq m$ (perché $m < n$ e $n - 1$ è il massimo dei numeri minori di $n$), cioè $n - 1 \in \mbN_m$.
        \item $n - 1 < n = \min(y)$, quindi $n - 1 \notin y$: deve essere $n - 1 \in x$.
        \item Per ipotesi allora $(n - 1) + 1 \in x$, cioè $n \in x$ (Proposizione~\ref{prop:successore}): ma $n \in y$, assurdo.
    \end{passi}
\end{dimostrazione}
\end{concetto}

\begin{spiegazione}
    È l'effetto domino: se cade la prima tessera ($m \in x$) e ogni tessera che cade fa cadere la successiva ($n \in x \Rightarrow n + 1 \in x$),
    cadono tutte le tessere da $m$ in poi.
    \begin{center}
    \begin{tikzpicture}
        \fill[black!12] (-0.5,-0.06) rectangle (8.9,0);
        \draw[fill=orange!35, draw=orange!80!black, rotate around={-58:(0.12,0)}] (-0.12,0) rectangle (0.12,1.3);
        \draw[fill=orange!35, draw=orange!80!black, rotate around={-32:(1.32,0)}] (1.08,0) rectangle (1.32,1.3);
        \foreach \i in {2,...,7} \draw[fill=primary!15, draw=primary!70!black] (\i*1.2-0.12,0) rectangle (\i*1.2+0.12,1.3);
        \foreach \i/\t in {0/m, 1/m+1, 2/m+2, 3/m+3, 4/m+4, 5/m+5, 6/m+6, 7/m+7} \node[font=\scriptsize] at (\i*1.2,-0.32) {$\t$};
        \draw[-{Stealth}, orange!80!black, line width=0.8pt] (0.5,1.45) to[bend left=30] (1.25,1.5);
        \draw[-{Stealth}, orange!80!black, line width=0.8pt] (1.75,1.5) to[bend left=30] (2.35,1.5);
        \node[font=\small, align=left, anchor=west] at (9.2,0.65) {base: $m \in x$\\ passo: $n \in x \Rightarrow n + 1 \in x$};
    \end{tikzpicture}
    \end{center}
    La dimostrazione formalizza un ragionamento per assurdo: se qualche tessera restasse in piedi, ce ne sarebbe una \textbf{prima}
    rimasta in piedi ($n = \min y$, esiste per il buon ordine); ma la tessera subito prima di lei ($n - 1$) è caduta, e cadendo l'avrebbe
    fatta cadere. È qui che servono le proprietà di $\mbN$: il minimo di ogni parte non vuota e il predecessore.
\end{spiegazione}

Nella pratica il principio si usa così.

\begin{concetto}
\begin{Corollario}{Dimostrazioni per induzione}{induzione-uso}
    \begin{prerequisiti}
        \dip{Teorema}{teo:induzione-prima}
    \end{prerequisiti}
    Sia $P(n)$ una proprietà dei numeri naturali e sia $m \in \mbN$. Supponiamo che
    \begin{itemize}
        \item[] \textbf{(passo base)} $P(m)$ sia vera;
        \item[] \textbf{(passo induttivo)} per ogni $n \geq m$, se $P(n)$ è vera (\textbf{ipotesi induttiva}), allora $P(n + 1)$ è vera.
    \end{itemize}
    Allora $P(n)$ è vera per ogni $n \geq m$.
\end{Corollario}
\begin{dimostrazione}
    Sia $x = \{n \in \mbN_m \mid P(n)\}$. Per il passo base $m \in x$, quindi $x \neq \varnothing$, e $\min(x) = m$ (tutti gli elementi
    di $x$ sono $\geq m$). Per il passo induttivo, $n \in x \rightarrow n + 1 \in x$. Per il Teorema~\ref{teo:induzione-prima}
    $x = \mbN_m$: ogni $n \geq m$ soddisfa $P(n)$.
\end{dimostrazione}
\end{concetto}

\begin{spiegazione}
    \begin{passi}
        \item Si dimostra il \textbf{passo base}: di solito un conto diretto per il primo valore.
        \item Si fissa un $n$ qualunque, si \textbf{suppone} vera $P(n)$ e si dimostra $P(n + 1)$, usando $P(n)$.
        \item Si conclude che $P$ vale da $m$ in poi.
    \end{passi}
    Attenzione: nel passo induttivo non si assume che $P$ valga per tutti gli $n$ (sarebbe la tesi!), ma solo per un $n$ fissato.
\end{spiegazione}

\begin{esempio}{La somma dei primi $n$ numeri}
    Dimostriamo che per ogni $n \in \mbN$: \quad $0 + 1 + \dots + (n - 1) = \dfrac{n(n - 1)}{2}$.
    \begin{passi}
        \item \textbf{Base} ($n = 0$): a sinistra la somma è vuota, cioè $0$; a destra $\frac{0 \cdot (-1)}{2} = 0$.
        \item \textbf{Passo}: supponiamo $0 + 1 + \dots + (n - 1) = \frac{n(n - 1)}{2}$. Aggiungendo $n$ a entrambi i membri:
              \[ 0 + 1 + \dots + (n - 1) + n = \frac{n(n - 1)}{2} + n = \frac{n(n - 1) + 2n}{2} = \frac{(n + 1)n}{2}, \]
              che è la formula per $n + 1$.
    \end{passi}
    Il disegno spiega anche \textbf{da dove viene} la formula: due «scale» da $0 + 1 + 2 + 3$ quadretti formano un rettangolo
    $4 \times 3$, quindi una scala ha $\frac{4 \cdot 3}{2} = 6$ quadretti.
    \begin{center}
    \begin{tikzpicture}[scale=0.45]
        \foreach \c in {1,2,3} \foreach \r in {1,...,\c} \draw[fill=primary!25, draw=primary!70!black] (\c,\r-1) rectangle (\c+1,\r);
        \foreach \c in {0,1,2} \foreach \r in {\c,...,2} \draw[fill=orange!30, draw=orange!80!black] (\c,\r) rectangle (\c+1,\r+1);
        \draw[decorate, decoration={brace, amplitude=4pt, mirror}] (0,-0.2) -- node[below=3pt, font=\small] {$4$} (4,-0.2);
        \draw[decorate, decoration={brace, amplitude=4pt}] (-0.2,0) -- node[left=3pt, font=\small] {$3$} (-0.2,3);
    \end{tikzpicture}
    \end{center}
\end{esempio}

A volte, per dimostrare $P(n)$, non basta sapere che vale $P(n - 1)$: serve che valga per \textbf{tutti} i numeri precedenti.

\begin{concetto}
\begin{Teorema}{Principio di induzione di seconda forma (induzione completa)}{induzione-seconda}
    \begin{prerequisiti}
        \dip{Definizione}{def:naturali}
        \dip{Definizione}{def:ben-ordinato}
    \end{prerequisiti}
    Sia $x$ una parte non vuota di $\mbN$ tale che, per ogni $n \in \mbN$, se tutti i $k$ con $\min(x) \leq k < n$ stanno in $x$,
    allora $n \in x$. Allora $x = \mbN_{\min(x)}$. In formule:
    \[ \forall x \in \mcP(\mbN) \setminus \{\varnothing\}\ \Big(\forall n \in \mbN\,\big(\forall k \in \mbN\,(\min(x) \leq k < n \rightarrow k \in x) \rightarrow n \in x\big)
       \rightarrow x = \mbN_{\min(x)}\Big). \]
\end{Teorema}
\begin{dimostrazione}
    Sia $m = \min(x)$; come prima $x \subseteq \mbN_m$. Per assurdo sia $y = \mbN_m \setminus x \neq \varnothing$ e sia $n = \min(y)$.
    Se $k$ è tale che $m \leq k < n$, allora $k \in \mbN_m$ e $k \notin y$ (perché $k < \min(y)$), quindi $k \in x$. Dunque tutti i $k$
    con $m \leq k < n$ stanno in $x$, e per ipotesi $n \in x$: ma $n \in y$, assurdo.
\end{dimostrazione}
\end{concetto}

\begin{spiegazione}
    È ancora l'effetto domino, ma con tessere più pesanti: per far cadere la tessera $n$ servono \textbf{tutte} quelle prima di lei.
    Il ragionamento per assurdo è lo stesso: la prima tessera rimasta in piedi avrebbe tutte le precedenti cadute, quindi cadrebbe
    anche lei. Nella pratica: si suppone $P(k)$ vera per \textbf{ogni} $k$ con $m \leq k < n$ e si dimostra $P(n)$. Quando il ragionamento
    usa uno dei $k$ precedenti ma non funziona per i primi valori, questi si controllano a mano come passi base.
\end{spiegazione}

\begin{esempio}{Francobolli da $4$ e da $5$}
    Dimostriamo che ogni $n \geq 12$ si scrive $n = 4a + 5b$ con $a, b \in \mbN$ (ogni affrancatura da almeno $12$ centesimi si fa
    con francobolli da $4$ e da $5$).
    \begin{passi}
        \item \textbf{Basi}: $12 = 4 \cdot 3$, \ $13 = 4 \cdot 2 + 5$, \ $14 = 4 + 5 \cdot 2$, \ $15 = 5 \cdot 3$.
        \item \textbf{Passo} ($n > 15$): supponiamo l'enunciato vero per ogni $k$ con $12 \leq k < n$. Allora $12 \leq n - 4 < n$, quindi
              $n - 4 = 4a + 5b$ e $n = 4(a + 1) + 5b$.
    \end{passi}
    \begin{center}
    \begin{tikzpicture}[x=0.8cm]
        \foreach \n in {12,...,23} {
            \pgfmathtruncatemacro{\r}{mod(\n,4)}
            \ifcase\r \def\col{classeA}\or \def\col{classeB}\or \def\col{classeC}\or \def\col{classeD}\fi
            \node[circle, fill=\col, draw=black!30, inner sep=0pt, minimum size=6.5mm, font=\small] (k\n) at (\n,0) {$\n$};
        }
        \foreach \n in {12,...,15} \node[font=\scriptsize] at (\n,-0.6) {base};
        \foreach \n in {12,13,14,15,16,17,18,19} {
            \pgfmathtruncatemacro{\m}{\n+4}
            \draw[-{Stealth}, line width=0.6pt, black!60] (k\n.north) to[bend left=40] (k\m.north);
        }
    \end{tikzpicture}
    \end{center}
    Ogni freccia aggiunge un francobollo da $4$: partendo dalle quattro basi si raggiungono tutti i numeri da $12$ in poi. Qui
    serve la seconda forma perché da $n$ si torna a $n - 4$, non a $n - 1$.
\end{esempio}

% ──────────────────────────────────────────────────────────────
\section{Potenze e multipli}\label{sec:potenze}

Il principio di induzione permette anche di \textbf{definire} oggetti «passo per passo»: si dice cosa succede in un primo
valore e come si passa da $n$ a $n + 1$. Si parla di \textbf{definizione ricorsiva}. Così si definiscono le potenze.

\dfn[potenze]{Potenze e multipli}{
    Sia $(s, *)$ un semigruppo e sia $x \in s$. Le \textbf{potenze} di $x$ sono definite ricorsivamente da
    \[ x^1 := x, \qquad x^{n + 1} := x^n * x \quad (n \geq 1). \]
    Se $(s, *, 1_s)$ è un monoide si pone anche $x^0 := 1_s$. Se $(s, *)$ è un gruppo e $n \geq 1$ si pone $x^{-n} := (x^{-1})^n$.
    In notazione additiva si parla di \textbf{multipli} e si scrive $nx$ invece di $x^n$: $1x = x$, $(n + 1)x = nx + x$, $0x = 0$,
    $(-n)x = n(-x)$.
}

\begin{concetto}
\begin{Teorema}{Regole delle potenze}{regole-potenze}
    \begin{prerequisiti}
        \dip{Definizione}{def:potenze}
        \dip{Corollario}{cor:induzione-uso}
        \dip{Definizione}{def:semigruppo}
    \end{prerequisiti}
    Sia $(s, *)$ un semigruppo e sia $x \in s$. Per ogni $m, n \geq 1$:
    \[ \text{(a)}\ \ x^m * x^n = x^{m + n} \qquad\qquad \text{(b)}\ \ (x^m)^n = x^{mn}. \]
    Se $s$ è un monoide, (a) e (b) valgono per ogni $m, n \geq 0$.
\end{Teorema}
\begin{dimostrazione}
    Fissiamo $m \geq 1$ e procediamo per induzione su $n$.
    \begin{caso}{(a) $x^m * x^n = x^{m + n}$}
        \textbf{Base} $n = 1$: $x^m * x^1 = x^m * x = x^{m + 1}$ per definizione. \textbf{Passo}: se $x^m * x^n = x^{m + n}$, allora, usando
        l'associatività,
        \[ x^m * x^{n + 1} = x^m * (x^n * x) = (x^m * x^n) * x = x^{m + n} * x = x^{m + n + 1}. \]
    \end{caso}
    \begin{caso}{(b) $(x^m)^n = x^{mn}$}
        \textbf{Base} $n = 1$: $(x^m)^1 = x^m$. \textbf{Passo}: se $(x^m)^n = x^{mn}$, allora per definizione e per (a)
        \[ (x^m)^{n + 1} = (x^m)^n * x^m = x^{mn} * x^m = x^{mn + m} = x^{m(n + 1)}. \]
    \end{caso}
    \begin{caso}{Esponenti nulli in un monoide}
        $x^m * x^0 = x^m * 1_s = x^m = x^{m + 0}$, \ $x^0 * x^n = x^n$, \ $(x^m)^0 = 1_s = x^0 = x^{m \cdot 0}$, \ $(x^0)^n = 1_s^n = 1_s = x^0$
        (l'ultima perché $1_s^{n + 1} = 1_s^n * 1_s = 1_s^n$, per induzione).
    \end{caso}
\end{dimostrazione}
\end{concetto}

\begin{spiegazione}
    Nel conto $x^2 * x^3 = (x * x) * (x * x * x) = x * x * x * x * x = x^5$ «togliamo le parentesi» grazie all'associatività e contiamo
    le $x$. Ma «e così via» non è una dimostrazione: l'induzione è il modo rigoroso di dirlo. Seguiamo il passo induttivo per $m = 2$:
    \begin{passi}
        \item $x^2 * x^1 = x^2 * x = x^3$: è la definizione di $x^3$ (base).
        \item $x^2 * x^2 = x^2 * (x^1 * x) = (x^2 * x^1) * x = x^3 * x = x^4$: si usa il passo precedente.
        \item $x^2 * x^3 = x^2 * (x^2 * x) = (x^2 * x^2) * x = x^4 * x = x^5$: e così via, ogni volta con il risultato del passo prima.
    \end{passi}
\end{spiegazione}

\begin{concetto}
\begin{Corollario}{Le potenze di uno stesso elemento commutano}{potenze-commutano}
    \begin{prerequisiti}
        \dip{Teorema}{teo:regole-potenze}
    \end{prerequisiti}
    In un semigruppo $(s, *)$, per ogni $x \in s$ e $m, n \geq 1$ (anche $\geq 0$ in un monoide): $x^m * x^n = x^n * x^m$.
\end{Corollario}
\begin{dimostrazione}
    $x^m * x^n = x^{m + n} = x^{n + m} = x^n * x^m$, per il Teorema~\ref{teo:regole-potenze} e la commutatività della somma in $\mbN$.
\end{dimostrazione}
\end{concetto}

\begin{spiegazione}
    Anche se l'operazione non è commutativa, $x^2 * x^3$ e $x^3 * x^2$ sono entrambi «cinque $x$ in fila»: le potenze di uno
    \textbf{stesso} elemento non si accorgono della mancanza di commutatività. Con elementi diversi invece $x * y$ e $y * x$ possono
    essere diversi, come nel monoide delle parole: $ab \neq ba$.
\end{spiegazione}

\begin{concetto}
\begin{Proposizione}{Potenze dell'inverso}{potenze-inverso}
    \begin{prerequisiti}
        \dip{Definizione}{def:potenze}
        \dip{Proposizione}{prop:unita-gruppo}
        \dip{Corollario}{cor:potenze-commutano}
        \dip{Corollario}{cor:induzione-uso}
    \end{prerequisiti}
    Sia $(g, *)$ un gruppo e $x \in g$. Per ogni $n \geq 1$: $(x^{-1})^n = (x^n)^{-1}$, cioè $x^{-n}$ è l'inverso di $x^n$.
\end{Proposizione}
\begin{dimostrazione}
    Per induzione su $n$. \textbf{Base}: $(x^{-1})^1 = x^{-1} = (x^1)^{-1}$. \textbf{Passo}: se $(x^{-1})^n = (x^n)^{-1}$, allora
    \[ (x^{-1})^{n + 1} = (x^{-1})^n * x^{-1} = (x^n)^{-1} * x^{-1} = (x * x^n)^{-1} = (x^{n + 1})^{-1}, \]
    dove abbiamo usato $(a * b)^{-1} = b^{-1} * a^{-1}$ (Proposizione~\ref{prop:unita-gruppo}) e $x * x^n = x^1 * x^n = x^{n + 1}$.
\end{dimostrazione}
\end{concetto}

\begin{spiegazione}
    Per «disfare» $x^3 = x * x * x$ bisogna togliere le tre $x$ una alla volta, partendo dall'ultima: $x^{-1} * x^{-1} * x^{-1}$.
    In $(\mbZ, +)$ questo dice che l'opposto di $3x$ è $3(-x)$; in $(\mbQ \setminus \{0\}, \cdot)$ che l'inverso di $2^3 = 8$ è
    $(\frac{1}{2})^3 = \frac{1}{8}$.
\end{spiegazione}

\info{Potenze con esponente intero}{
    In un gruppo le regole $x^m * x^n = x^{m + n}$ e $(x^m)^n = x^{mn}$ valgono per \textbf{tutti} gli $m, n \in \mbZ$: si dimostra
    distinguendo i casi secondo i segni di $m$ e $n$ e usando la Proposizione~\ref{prop:potenze-inverso}. Ad esempio
    $x^5 * x^{-2} = x * x * x * x * x * x^{-1} * x^{-1} = x^3$.
}

Le potenze descrivono le sottostrutture generate da un solo elemento.

\begin{concetto}
\begin{Proposizione}{Sottostrutture cicliche}{cicliche}
    \begin{prerequisiti}
        \dip{Teorema}{teo:sottomonoidi-generati}
        \dip{Teorema}{teo:sottogruppi-generati}
        \dip{Definizione}{def:potenze}
        \dip{Teorema}{teo:regole-potenze}
    \end{prerequisiti}
    Sia $x$ un elemento di un monoide $(m, *, 1_m)$. Il sottomonoide generato da $x$ è $\langle x \rangle = \{x^n \mid n \in \mbN\}$.
    Se $m$ è un gruppo, il sottogruppo generato da $x$ è $\langle x \rangle = \{x^n \mid n \in \mbZ\}$.
\end{Proposizione}
\begin{dimostrazione}
    Per il Teorema~\ref{teo:sottomonoidi-generati} con $t = \{x\}$, $\langle x \rangle$ è formato da $1_m = x^0$ e dai prodotti
    $x * x * \dots * x$ di $n \geq 1$ fattori, che sono le potenze $x^n$. Nel caso di un gruppo, per il
    Teorema~\ref{teo:sottogruppi-generati} gli elementi di $\langle x \rangle$ sono i prodotti di fattori $x$ e $x^{-1}$; per le regole delle
    potenze con esponente intero ognuno di essi è una potenza $x^k$, con $k$ uguale al numero di fattori $x$ meno il numero di fattori
    $x^{-1}$; e ogni $x^k$ è un prodotto di questo tipo.
\end{dimostrazione}
\end{concetto}

\begin{spiegazione}
    Partendo da un solo elemento e potendo solo moltiplicarlo per sé stesso (e, in un gruppo, per il suo inverso) si ottengono
    solo le sue potenze. Ad esempio in $(\mbZ, +)$, in notazione additiva, $\langle 1 \rangle = \{n \cdot 1 \mid n \in \mbZ\} = \mbZ$ (il
    gruppo è \textbf{ciclico}) e $\langle m \rangle = \{nm \mid n \in \mbZ\} = m\mbZ$, l'insieme dei multipli di $m$; in $(\mbN, +, 0)$
    il sottomonoide generato da $2$ è $\{0, 2, 4, \dots\} = 2\mbN$.
\end{spiegazione}

Negli anelli i multipli permettono di confrontare la moltiplicazione per un intero con la moltiplicazione per un elemento
dell'anello (Definizione~\ref{def:notazioni-anello}).

\begin{concetto}
\begin{Proposizione}{Moltiplicazione per un multiplo dell'unità}{multipli-unita}
    \begin{prerequisiti}
        \dip{Definizione}{def:notazioni-anello}
        \dip{Lemma}{lem:regole-anello}
        \dip{Corollario}{cor:induzione-uso}
    \end{prerequisiti}
    Sia $(a, +, \cdot)$ un anello unitario. Per ogni $x \in a$ e $n \in \mbZ$: \quad $(n 1_a) \cdot x = nx = x \cdot (n 1_a)$.
\end{Proposizione}
\begin{dimostrazione}
    \begin{caso}{Passo 1: $(ky)x = k(yx)$ e $x(ky) = k(xy)$ per ogni $k \geq 1$ e $y \in a$}
        Per induzione su $k$. Base: $(1y)x = yx = 1(yx)$. Passo: per la distributività,
        $((k + 1)y)x = (ky + y)x = (ky)x + yx = k(yx) + yx = (k + 1)(yx)$. A destra è identico.
    \end{caso}
    \begin{caso}{Passo 2: i tre casi di $n$}
        Se $n \geq 1$, per il passo 1 con $y = 1_a$: $(n1_a)x = n(1_a x) = nx$. Se $n = 0$: $(0 \cdot 1_a)x = 0_a x = 0_a = 0x$ (Lemma~\ref{lem:regole-anello}).
        Se $n = -k$ con $k \geq 1$: per definizione $n1_a = k(-1_a)$, quindi per il passo 1 e il Lemma~\ref{lem:regole-anello}
        $(n1_a)x = k((-1_a)x) = k(-x) = nx$. A destra è identico.
    \end{caso}
\end{dimostrazione}
\end{concetto}

\begin{spiegazione}
    Moltiplicare per «$1 + 1 + 1$» equivale a sommare tre volte: $(3 \cdot 1_a)x = (1_a + 1_a + 1_a)x = x + x + x = 3x$. La proposizione
    dice che questo vale per ogni intero, anche negativo, e sia a sinistra sia a destra, anche in un anello non commutativo.
    In $\mbZ$ sembra ovvio; in un anello come $(\mcP(s), \Delta, \cap)$ dice ad esempio che $(2 \cdot s) \cap x = 2x = x \,\Delta\, x = \varnothing$.
\end{spiegazione}
